% ============================================================================
% preamble.tex -- Self-reinforced geometric absorption
% Engine: pdflatex (see latexmkrc).  Bibliography: biblatex + biber, style=phys
% ============================================================================

% ---- Layout & typography ---------------------------------------------------
\usepackage[a4paper,margin=2.5cm,headheight=14pt]{geometry}
\usepackage[T1]{fontenc}
\usepackage[utf8]{inputenc}
\usepackage{lmodern}
\usepackage{microtype}
\usepackage{csquotes}

% ---- Mathematics -----------------------------------------------------------
\usepackage{amsmath}
\usepackage{amssymb}
\usepackage{amsthm}
\usepackage{mathtools}
\usepackage{bm}
\usepackage{mathrsfs}   % \mathscr, used for the Hilbert space
\usepackage{tensor}

% ---- Graphics & tables -----------------------------------------------------
\usepackage{graphicx}
\usepackage{float}
\usepackage[font=small,labelfont=bf,labelsep=period]{caption}
\usepackage{subcaption}
\graphicspath{{figures/}}
\usepackage{ifthen}
\usepackage{booktabs}
\usepackage{array}
\usepackage{longtable}
\usepackage{tabularx}
\usepackage{multirow}
\usepackage[group-separator={,},group-minimum-digits=4,detect-all]{siunitx}
% siunitx knows no astronomical units; the parsec is needed for H_0.
\DeclareSIUnit{\parsec}{pc}
\newcolumntype{L}[1]{>{\raggedright\arraybackslash}p{#1}}
\newcolumntype{C}[1]{>{\centering\arraybackslash}p{#1}}
\newcolumntype{Y}{>{\raggedright\arraybackslash}X}

% ---- Lists, spacing, floats ------------------------------------------------
\usepackage{enumitem}
\usepackage{xspace}
\usepackage{etoolbox}
\usepackage[section]{placeins}
\usepackage{needspace}
\usepackage{xcolor}

% ---- Links (load before cleveref) ------------------------------------------
\usepackage[colorlinks=true,linkcolor=blue!55!black,citecolor=green!45!black,
            urlcolor=black!65,bookmarks=true,bookmarksnumbered=true]{hyperref}
\usepackage[nameinlink,capitalise]{cleveref}

% ---- Bibliography ----------------------------------------------------------
\usepackage[backend=biber,style=phys,biblabel=brackets,sorting=none,
            eprint=true,doi=true,url=false,isbn=false,
            maxbibnames=5,minbibnames=3,giveninits=true]{biblatex}
\addbibresource{references.bib}

% ============================================================================
% Draft-figure guard.  In draft mode a missing figure prints a visible box;
% in final mode a missing figure is a hard error (never silently dropped).
% ============================================================================
\newif\ifdraftfigs
\draftfigstrue   % <-- set \draftfigsfalse before submission
\let\origincludegraphics\includegraphics
\renewcommand{\includegraphics}[2][]{%
  \IfFileExists{figures/#2}%
    {\origincludegraphics[#1]{#2}}%
    {\ifdraftfigs
       \fbox{\parbox{0.6\linewidth}{\centering\color{red}\ttfamily
             MISSING FIGURE\\ figures/#2}}%
     \else
       \GenericError{}{Missing figure: figures/#2}{}{Build aborted: final mode requires all figures.}%
     \fi}}

% ============================================================================
% Theorem environments (shared counter, numbered by section)
% ============================================================================
\theoremstyle{plain}
\newtheorem{theorem}{Theorem}[section]
\newtheorem{proposition}[theorem]{Proposition}
\newtheorem{lemma}[theorem]{Lemma}
\newtheorem{corollary}[theorem]{Corollary}
\theoremstyle{definition}
\newtheorem{definition}[theorem]{Definition}
\newtheorem{assumption}[theorem]{Assumption}
\newtheorem{example}[theorem]{Example}
\newtheorem{conjecture}[theorem]{Conjecture}
\theoremstyle{remark}
\newtheorem{remark}[theorem]{Remark}
\newtheorem{notation}[theorem]{Notation}

\crefname{theorem}{Theorem}{Theorems}
\crefname{proposition}{Proposition}{Propositions}
\crefname{lemma}{Lemma}{Lemmas}
\crefname{corollary}{Corollary}{Corollaries}
\crefname{definition}{Definition}{Definitions}
\crefname{assumption}{Assumption}{Assumptions}
\crefname{conjecture}{Conjecture}{Conjectures}
\crefname{remark}{Remark}{Remarks}
\crefname{equation}{Eq.}{Eqs.}
\crefname{figure}{Fig.}{Figs.}
\crefname{table}{Table}{Tables}
\crefname{section}{Sec.}{Secs.}
\crefname{appendix}{Appendix}{Appendices}

\numberwithin{equation}{section}

% ============================================================================
% Claims ladder.  Every displayed key result is wrapped in `keyresult`, whose
% level is machine-extracted by scripts/check_claims.py:
%   I = identity or standard result   C = conditional theorem of the model
%   P = phenomenological / fitted     S = speculative (Outlook only)
% Usage:  \begin{keyresult}{C}{eq:jerk-bound} ... \end{keyresult}
% ============================================================================
% The level is printed in the margin beside the result.  It used to be
% discarded, so a reader had to match labels against a table in an appendix to
% learn whether a displayed result was an identity, a theorem or a conjecture;
% a ladder invisible at the point of the claim does not do the job the
% introduction claims for it.  The argument is also extracted from the source
% by scripts/check_claims.py, which greps \begin{keyresult}{LEVEL}{label}.
\newcommand{\claimlevel}[1]{%
  \marginpar{\raggedright\scriptsize\sffamily\textbf{#1}\\[-1pt]%
    \tiny\MakeLowercase{\claimlevelname{#1}}}}
\newcommand{\claimlevelname}[1]{%
  \ifthenelse{\equal{#1}{A}}{Postulate}{%
  \ifthenelse{\equal{#1}{I}}{Identity}{%
  \ifthenelse{\equal{#1}{C}}{Theorem}{%
  \ifthenelse{\equal{#1}{N}}{Numerical}{%
  \ifthenelse{\equal{#1}{P}}{Fitted}{%
  \ifthenelse{\equal{#1}{S}}{Conjecture}{}}}}}}}
\newenvironment{keyresult}[2]%
  {\claimlevel{#1}\begin{equation}\label{#2}\begin{aligned}}%
  {\end{aligned}\end{equation}}

% ---- Draft helpers (removed before submission) ------------------------------
\newcommand{\TODO}[1]{\textcolor{red}{\textbf{[TODO: #1]}}}
\newcommand{\CITE}{\textcolor{orange}{\textbf{[CITE]}}}
